\subsection{An exact commutator with unbounded quotient norm}
The following calculation records why a shrinking parity gap requires
an estimate beyond strong convergence and compactness.
On two copies of $\ell^2(\mathbb N)$ let
$D=\operatorname{diag}(2^{-j})_{j\ge1}$ and
$C=\left(\begin{smallmatrix}0&D\\D&0\end{smallmatrix}\right)$.
For $n\ge1$ put $\tau=2^{-n}$, $\kappa=\tau^4$, $r=1-\kappa$, and
$v_{n+k}=2^{-(k+1)/2}$ for $k\ge0$, with all earlier coordinates zero.
Thus $\|v\|=1$ and $v\notin\operatorname{ran}D$. Define
\[
 H=\tau^4D^2,\quad
 a=\langle(I-H)^{-1}v,v\rangle^{-1},\quad
 b=\langle(rI-H)^{-1}Dv,Dv\rangle^{-1},
\]
\[
 M_+=H+a\,v\otimes v^*,\qquad
 M_-=H+b\,Dv\otimes(Dv)^*,\qquad S_\pm=I-M_\pm.
\]
The top eigenvalues of $M_+,M_-$ are respectively $1,r$, both simple.
Indeed, for $t>\|H\|$ and $c=\langle(tI-H)^{-1}u,u\rangle^{-1}$,
\[
 tI-H-c\,u\otimes u^*
 =(tI-H)^{1/2}(I-cw\otimes w^*)(tI-H)^{1/2},\quad
 w=(tI-H)^{-1/2}u,
\]
and $c\|w\|^2=1$. This proves positivity and the one-dimensional
kernel. Each $M_\pm$ is compact and strictly positive on nonzero
vectors because $H$ has positive diagonal entries.

Let $e,f$ be the unit normalizations of $(I-H)^{-1}v$ and
$(rI-H)^{-1}Dv$. The tail supports escape to infinity,
$a\le1$, and $b\|Dv\|^2\le r$, so $M_\pm\to0$ strongly.
Finite-rank approximation of $D$ yields
$\|M_-D\|_{\mathcal S_2}+\|DM_+\|_{\mathcal S_2}\to0$.
Nevertheless the quotient operator below becomes unbounded in norm.

The exact commutator is
\[
 DM_+-M_-D=Dv\otimes\ell^0,\qquad
 \ell^0(h)=a\langle h,v\rangle-b\langle h,D^2v\rangle.
\]
Evaluating at $e$ gives $S_-De=Dv\,\ell^0(e)$.
Since $S_-\ge\kappa I$ and $D$ is injective, $\ell^0(e)\ne0$.
Set $\ell=\ell^0/\ell^0(e)$ on the even copy and zero on the odd
copy. Then
\[
 DM_+-M_-D=S_-De\otimes\ell,\qquad
 D(I-e\otimes\ell)=S_-^{-1}DS_+.
\]
The quotient of $C(I-e\otimes\ell)$ by $\mathbb Ce$, with metric
$I-M$, is therefore represented isometrically by
\[
 B_n=\begin{pmatrix}0&S_+^{1/2}DS_-^{-1/2}\\
 S_-^{-1/2}DS_+^{1/2}&0\end{pmatrix}
\]
off its zero vector $e$. At each $n$ it is compact and selfadjoint.
It is injective on the quotient: the lower-left block has kernel
$\mathbb Ce$; the kernel of its adjoint requires
$DS_-^{-1/2}h\in\mathbb Ce$. Since
$e_j$ is proportional to $v_j/(1-\tau^4 2^{-2j})$, it does not belong
to $\operatorname{ran}D$, and hence $h=0$.

For the norm, put $f^0=Dv/\|Dv\|$. Exact geometric sums give
\[
 \|Dv\|^2=\frac47\tau^2,\quad
 \|D^2v\|^2=\frac{16}{31}\tau^4,\quad
 \langle Df^0,(I-v\otimes v^*)Df^0\rangle
       =\frac{72}{217}\tau^2.
\]
On the common tail, $\|D\|=\tau$ and $\|H\|=\eta=\tau^6$.
Since $r\ge15/16$ and $\eta\le1/64$, the resolvent geometric series
and normalization give
$\|f-f^0\|\le2(\eta/r)/(1-\eta/r)\le3\eta$.
Also $1-\eta\le a\le1$, so
$\|S_+-(I-v\otimes v^*)\|\le2\eta$ on that tail.
It follows that
\[
 \langle Df,S_+Df\rangle
 \ge\frac{72}{217}\tau^2-8\tau^8
 \ge\frac{359}{1736}\tau^2>\frac15\tau^2.
\]
Since $S_-f=\kappa f$, one obtains
\[
 \|B_n\|^2\ge\|B_n(0,f)\|^2
       =\kappa^{-1}\langle Df,S_+Df\rangle
       \ge\frac1{5\tau^2}\longrightarrow\infty.
\]
The construction realizes the same rank-one commutator equation and
compact-product limits appearing in the parity comparison. Its
quantified divergence shows that those displayed data alone do not
bound the quotient norm. Applying the parity method to the Weil
operator still requires control of its gap or of the weighted product
in \eqref{eq:parity-error}.
